\chapter{User Manual}
\label{app:user-manual}

\section{Running the Tool Locally}
From the repository root, start a local server:
\begin{verbatim}
python -m http.server 5500
\end{verbatim}
Open:
\begin{verbatim}
http://localhost:5500/
\end{verbatim}

\section{Basic Usage}
\begin{itemize}
  \item Select an algorithm (Token Ring or RA) and choose the number of processes.
  \item Use \textbf{Request CS} to issue a request.
  \item Use \textbf{Step} to progress deterministically.
  \item Use \textbf{Run/Pause} for automated playback.
  \item Observe the process table, message queue table, trace, safety label, and preview canvas.
\end{itemize}

\section{Fault Injection}
\begin{itemize}
  \item Token Ring: use \textbf{Drop token}, \textbf{Regenerate token}, \textbf{Crash}, and \textbf{Recover}.
  \item RA: use \textbf{Drop next in-flight msg} and \textbf{Arm drop-next-send}.
\end{itemize}

\section{Scripted Demos}
Use the scripted demo controls to load deterministic scenarios such as crash/recover and RA tie-break demonstrations.

\section{Evidence Export}
Use the export controls to generate:
\begin{itemize}
  \item state JSON,
  \item trace TXT,
  \item preview PNG.
\end{itemize}
These artefacts support reproducible evidence collection.